% =====================================================================
\chapter{Modulo 13 -- Alberi Binari e Alberi Binari di Ricerca (BST)}
% =====================================================================

\tikzset{
  albero/.style={every node/.style={circle,draw,fill=heapcol,font=\small,minimum size=7mm,inner sep=1pt},
                 level distance=1.05cm, edge from parent/.style={draw,thick}},
  evid/.style={fill=green!30!sfondo},
  tolto/.style={fill=red!25!sfondo},
}

% ---------------------------------------------------------------------
\section{Teoria \& Analogia Chiave}
% ---------------------------------------------------------------------

\subsection*{Dalle liste agli alberi}

In una lista ogni nodo ha \textbf{un solo} successivo. In un \textbf{albero} ogni nodo
può avere \textbf{più} successori, chiamati \emph{figli}. È la struttura naturale per
tutto ciò che è gerarchico: le cartelle del computer, l'organigramma di un'azienda,
l'albero genealogico, i capitoli e paragrafi di un libro.

\begin{analogia}[: l'albero genealogico (capovolto)]
Gli alberi dell'informatica si disegnano a testa in giù: la \textbf{radice} in alto,
le \textbf{foglie} in basso. Come in un albero genealogico, ogni nodo ha un
\textbf{padre} (tranne la radice) e zero o più \textbf{figli}; nodi con lo stesso padre
sono \textbf{fratelli}; risalendo si incontrano gli \textbf{antenati}, scendendo i
\textbf{discendenti}.
\end{analogia}

Il vocabolario delle slide (da sapere per i quiz):
\begin{itemize}
  \item un albero è un insieme di nodi collegati da archi in cui ogni nodo ha
        \textbf{al più un arco entrante} (un solo padre); il nodo senza padre è la
        \textbf{radice}, ed è unica;
  \item una \textbf{foglia} è un nodo senza figli; gli altri sono \textbf{nodi interni};
  \item il \textbf{livello} (o profondità) di un nodo è la sua distanza dalla radice: la
        radice è a livello 0, i suoi figli a livello 1, e così via;
  \item la \textbf{profondità} (o \textbf{altezza}) dell'albero è la lunghezza del
        cammino più lungo dalla radice a una foglia;
  \item ogni nodo è la radice di un \textbf{sottoalbero}, formato da lui e da tutti i suoi
        discendenti.
\end{itemize}

\textbf{Definizione ricorsiva} (la più importante): un albero è
\begin{itemize}
  \item \textbf{vuoto},
  \item oppure è formato da un nodo radice da cui discendono zero o più
        \textbf{sottoalberi}, che sono a loro volta alberi.
\end{itemize}
Un albero ``contiene'' alberi più piccoli: per questo quasi tutti gli algoritmi sugli
alberi sono \textbf{ricorsivi}, con il caso base sull'albero vuoto (Modulo 8). Riguarda le
``tre domande'' del Modulo 8: qui la risposta è quasi sempre ``faccio qualcosa con la
radice e chiedo il resto ai due sottoalberi''.

\subsection*{Alberi binari}

Un albero è \textbf{binario} se ogni nodo ha \textbf{al più due figli}, distinti in
\textbf{sinistro} e \textbf{destro}. In C++:
\begin{lstlisting}
template <typename T>
struct Nodo {
    T val;
    Nodo<T>* left = nullptr;     // figlio sinistro
    Nodo<T>* right = nullptr;    // figlio destro
    Nodo<T>* padre = nullptr;    // facoltativo: serve per cancellazione e successore
    Nodo(T v) : val(v) {}
};
\end{lstlisting}
È come il nodo di una lista doppia, ma i due puntatori vanno ``in giù'' verso due figli
invece che avanti e indietro.

Due numeri utili (dalle slide): al livello $k$ ci sono \textbf{al più} $2^k$ nodi, e un
albero binario di altezza $h$ ha al più $2^{h+1} - 1$ nodi. Rovesciando: con $n$ nodi,
un albero ben ``riempito'' (\textbf{bilanciato}) ha altezza circa $\log_2 n$; un albero
\textbf{degenere}, in cui ogni nodo ha un solo figlio, è in pratica una lista e ha
altezza $n - 1$.

\subsection*{Le tre visite}

\textbf{Visitare} un albero significa passare da tutti i nodi, una volta ciascuno.
Per un albero binario ci sono tre ordini classici, che differiscono solo per
\textbf{quando} si visita la radice rispetto ai due sottoalberi:
\begin{center}
\begin{tabular}{lll}
\toprule
\textbf{Visita} & \textbf{Ordine} & \textbf{Uso tipico} \\
\midrule
pre-order  & radice, sinistro, destro & copiare l'albero, stamparne la struttura \\
in-order   & sinistro, radice, destro & in un BST: valori \textbf{in ordine crescente} \\
post-order & sinistro, destro, radice & distruggere l'albero (prima i figli, poi il padre) \\
\bottomrule
\end{tabular}
\end{center}
\begin{lstlisting}
void inOrder(Nodo<T>* p) {
    if (p != nullptr) {              // caso base: albero vuoto, non faccio niente
        inOrder(p->left);            // 1. tutto il sottoalbero sinistro
        cout << p->val << " ";       // 2. la radice
        inOrder(p->right);           // 3. tutto il sottoalbero destro
    }
}
\end{lstlisting}
Spostando la riga del \cc|cout| prima o dopo le due chiamate si ottengono pre-order e
post-order: è lo stesso ``prima o dopo la chiamata'' del Modulo 8. Ogni visita tocca ogni
nodo una volta: costo $\Theta(n)$.

C'è anche una quarta visita, \textbf{per livelli} (in ampiezza): prima la radice, poi
tutti i nodi di livello 1, poi quelli di livello 2\ldots{} Non è ricorsiva: usa una
\textbf{coda} (Modulo 12; Esercizio~13.6). È la stessa idea della visita BFS dei grafi
(Modulo 14).

\subsection*{L'albero binario di ricerca (BST)}

\begin{definizione}[: proprietà del BST]
Un \textbf{albero binario di ricerca} (\emph{Binary Search Tree}) è un albero binario
in cui, per \textbf{ogni} nodo $x$:
\begin{itemize}
  \item tutti i valori nel sottoalbero \textbf{sinistro} di $x$ sono \textbf{minori} di
        $x$;
  \item tutti i valori nel sottoalbero \textbf{destro} di $x$ sono \textbf{maggiori o
        uguali} a $x$.
\end{itemize}
\end{definizione}
Attenzione: la regola vale per \emph{tutto} il sottoalbero, non solo per i figli diretti
(Esercizio~13.8).

\begin{analogia}[: il gioco ``più alto o più basso'']
Penso un numero da 1 a 100 e tu provi a indovinarlo. Dici 50: ``più basso''. Allora
scarti tutti i numeri da 50 in su e provi 25\ldots{} Un BST è questo gioco già
preparato: in ogni nodo chiedi ``più piccolo o più grande?'' e scendi a sinistra o a
destra, scartando ogni volta un intero sottoalbero. È la ricerca binaria del Modulo 7, ma
su una struttura dinamica in cui si può anche inserire e cancellare.
\end{analogia}

\textbf{Conseguenza fondamentale}: la visita \textbf{in-order} di un BST restituisce i
valori \textbf{in ordine crescente}. Inserire $n$ valori in un BST e fare una visita
in-order è un modo di ordinarli.

\subsection*{Ricerca, minimo e massimo}

\textbf{Ricerca} di un valore: parti dalla radice; se il valore è quello del nodo hai
finito; se è minore scendi a sinistra, altrimenti a destra; se arrivi a \cc|nullptr| il
valore non c'è.
\begin{lstlisting}
Nodo<T>* search(T elemento) const {
    Nodo<T>* x = root;
    while (x != nullptr && !(elemento == x->val))
        x = (elemento < x->val) ? x->left : x->right;
    return x;                    // nullptr se non trovato
}
\end{lstlisting}
\textbf{Minimo}: scendi sempre a sinistra finché puoi. \textbf{Massimo}: sempre a destra.

Tutte queste operazioni seguono un solo cammino dalla radice verso il basso: costano
$O(h)$, dove $h$ è l'altezza. Su un albero bilanciato $h \approx \log_2 n$ e la ricerca è
velocissima (1000 nodi: 10 confronti). Su un albero degenere $h = n - 1$ e la ricerca
diventa lineare (1000 confronti). Succede inserendo i valori \textbf{già ordinati}: ogni
nuovo nodo finisce sempre a destra (Esercizio~13.5).

\subsection*{Inserimento}

\begin{passopasso}[: inserimento in un BST (slide)]
I nuovi elementi vengono \textbf{sempre inseriti come foglie}.
\begin{enumerate}
  \item Parti dalla radice con un puntatore \cc|x|, e tieni in \cc|y| il padre di \cc|x|.
  \item Finché \cc|x| non è \cc|nullptr|: \cc|y = x|, poi scendi a sinistra se il
        valore è minore di \cc|x->val|, altrimenti a destra.
  \item Ora \cc|x| è il posto vuoto in cui va il nuovo nodo e \cc|y| è il suo futuro
        padre. Crea il nodo e agganciagli \cc|y| come padre.
  \item Se \cc|y| è \cc|nullptr| l'albero era vuoto: il nodo diventa la radice.
        Altrimenti diventa figlio sinistro o destro di \cc|y|, a seconda del confronto.
\end{enumerate}
\end{passopasso}
Esiste anche una versione ricorsiva molto compatta (usata negli esercizi):
\begin{lstlisting}
Nodo* inserisci(Nodo* r, int v) {
    if (r == nullptr) return new Nodo(v);          // posto vuoto: il nodo nasce qui
    if (v < r->val) r->left = inserisci(r->left, v);
    else r->right = inserisci(r->right, v);
    return r;                                      // la radice di questo sottoalbero non cambia
}
radice = inserisci(radice, 42);
\end{lstlisting}

\subsection*{Cancellazione: i tre casi}

È l'operazione più delicata, e una domanda d'orale classica. Per cancellare il nodo
\cc|z| bisogna mantenere la proprietà del BST. Tre casi:
\begin{enumerate}
  \item \textbf{\cc|z| è una foglia}: lo si toglie, e il puntatore del padre diventa
        \cc|nullptr|.
  \item \textbf{\cc|z| ha un solo figlio}: il figlio (con tutto il suo sottoalbero) prende
        il posto di \cc|z|. ``Il nonno adotta il nipote''.
  \item \textbf{\cc|z| ha due figli}: si cerca il suo \textbf{successore} \cc|y|, cioè
        il \textbf{minimo del sottoalbero destro}: è il valore immediatamente più grande di
        \cc|z|, e ha al più un figlio (quello destro: se avesse un figlio sinistro, quello
        sarebbe più piccolo). \cc|y| prende il posto di \cc|z|: così a sinistra restano
        tutti valori minori e a destra tutti valori maggiori. Se \cc|y| non è figlio
        diretto di \cc|z|, prima lo si stacca dal suo posto mettendo al suo posto il suo
        figlio destro.
\end{enumerate}
Le slide implementano tutto con una procedura di supporto, \textbf{\cc|Trapianta(u, v)|},
che sostituisce il sottoalbero di radice \cc|u| con quello di radice \cc|v|
aggiornando il padre di \cc|u|: con questa i casi 1 e 2 diventano una sola riga, e il caso
3 qualche riga in più (lo trovi nel pattern di codice). Costo: $O(h)$.

\textbf{Successore di un nodo} (slide): se il nodo ha un sottoalbero destro, è il minimo
di quel sottoalbero; altrimenti si risale verso la radice finché si arriva a un nodo
\emph{da sinistra}: quel padre è il successore (Esercizio~13.10).

\subsection*{Distruggere un albero}

Il distruttore deve liberare ogni nodo, e un nodo va liberato \textbf{dopo} i suoi figli
(altrimenti i puntatori ai figli si perdono): è una visita \textbf{post-order}.
\begin{lstlisting}
void distruggi(Nodo<T>* p) {
    if (p) { distruggi(p->left); distruggi(p->right); delete p; }
}
\end{lstlisting}

\begin{inparole}
\begin{itemize}
  \item Albero = radice + sottoalberi, che sono alberi: algoritmi ricorsivi, caso base
        sull'albero vuoto.
  \item Visite: pre (radice prima), in (radice in mezzo), post (radice dopo). In-order di
        un BST = valori ordinati.
  \item BST: a sinistra i minori, a destra i maggiori o uguali, in \emph{tutto} il sottoalbero.
  \item Ricerca, inserimento, cancellazione: $O(h)$. Bilanciato $h \approx \log n$,
        degenere $h \approx n$.
  \item Inserimento: sempre come nuova foglia.
  \item Cancellazione: foglia (si toglie), un figlio (il figlio sale), due figli (sale il
        successore, cioè il minimo del sottoalbero destro).
  \item Distruzione: in post-order.
\end{itemize}
\end{inparole}

% ---------------------------------------------------------------------
\section{Rappresentazione Visiva della Memoria (TikZ)}
% ---------------------------------------------------------------------

\textbf{Il BST del pattern} (prima delle cancellazioni), con i nomi delle parti.
\begin{center}
\begin{tikzpicture}[albero,level 1/.style={sibling distance=4.2cm},level 2/.style={sibling distance=2.1cm},level 3/.style={sibling distance=1.1cm}]
  \node {50}
    child {node {30}
      child {node {20}}
      child {node {40} child {node {35}} child {node {45}}}}
    child {node {70}
      child {node {60} child[missing] child {node {65}}}
      child {node {80}}};
  \node[draw=none,fill=none,font=\footnotesize\itshape,text=black!70,anchor=west] at (1.2,0.1) {radice (livello 0)};
  \node[draw=none,fill=none,font=\footnotesize\itshape,text=black!70,anchor=west] at (4.6,-2.1) {foglie: 20, 35, 45, 65, 80};
  \node[draw=none,fill=none,font=\footnotesize\itshape,text=black!70,anchor=west] at (4.6,-2.6) {altezza = 3};
\end{tikzpicture}
\end{center}

\needspace{6cm}
\textbf{Inserire 42}: il cammino dalla radice (in verde) e la nuova foglia.
\begin{center}
\begin{tikzpicture}[albero,level 1/.style={sibling distance=4.2cm},level 2/.style={sibling distance=2.1cm},level 3/.style={sibling distance=1.1cm}]
  \node[evid] {50}
    child {node[evid] {30}
      child {node {20}}
      child {node[evid] {40} child {node {35}} child {node[evid] {45} child {node[fill=green!55!sfondo] {42}} child[missing]}}}
    child {node {70}
      child {node {60}}
      child {node {80}}};
  \node[draw=none,fill=none,font=\footnotesize\itshape,text=black!70,anchor=west,text width=4.6cm] at (4.6,-3.0) {42 < 50: sinistra; 42 > 30: destra; 42 > 40: destra; 42 < 45: sinistra $\to$ posto vuoto};
\end{tikzpicture}
\end{center}

\needspace{6cm}
\textbf{Cancellazione con due figli}: si toglie 30. Il successore è il minimo del
sottoalbero destro di 30, cioè 35, che prende il suo posto.
\begin{center}
\begin{tikzpicture}[albero,level 1/.style={sibling distance=2.2cm},level 2/.style={sibling distance=1.1cm}]
  \node[tolto] at (0,0) {30}
    child {node {20}}
    child {node {40} child {node[evid] {35}} child {node {45}}};
  \node[draw=none,fill=none,font=\Large] at (2.8,-1) {$\Rightarrow$};
  \node[evid] at (5.4,0) {35}
    child {node {20}}
    child {node {40} child[missing] child {node {45}}};
  \node[draw=none,fill=none,font=\footnotesize\itshape,text=black!70,text width=5cm,anchor=west] at (7.2,-1) {a sinistra di 35 restano tutti valori minori, a destra tutti maggiori: la proprietà del BST è salva};
\end{tikzpicture}
\end{center}

\needspace{6cm}
\textbf{Bilanciato contro degenere}: gli stessi 7 valori inseriti in ordine diverso.
\begin{center}
\begin{tikzpicture}[albero,level 1/.style={sibling distance=2.2cm},level 2/.style={sibling distance=1.1cm}]
  \node at (0,0) {4}
    child {node {2} child {node {1}} child {node {3}}}
    child {node {6} child {node {5}} child {node {7}}};
  \node[draw=none,fill=none,font=\footnotesize\itshape,text=black!70] at (0,-3.0) {inseriti 4 2 6 1 3 5 7: altezza 2};
  \begin{scope}[xshift=5cm,level distance=0.62cm,level 1/.style={sibling distance=0.9cm},level 2/.style={sibling distance=0.9cm},level 3/.style={sibling distance=0.9cm},level 4/.style={sibling distance=0.9cm},level 5/.style={sibling distance=0.9cm},level 6/.style={sibling distance=0.9cm},every node/.style={circle,draw,fill=heapcol,font=\scriptsize,minimum size=5mm,inner sep=0pt}]
    \node {1} child[missing] child {node {2} child[missing] child {node {3} child[missing] child {node {4} child[missing] child {node {5} child[missing] child {node {6} child[missing] child {node {7}}}}}}};
  \end{scope}
  \node[draw=none,fill=none,font=\footnotesize\itshape,text=black!70,anchor=west] at (7.4,-1.0) {inseriti 1 2 3 4 5 6 7:};
  \node[draw=none,fill=none,font=\footnotesize\itshape,text=black!70,anchor=west] at (7.4,-1.5) {altezza 6 (una lista!)};
\end{tikzpicture}
\end{center}

% ---------------------------------------------------------------------
\section{Pattern di Codice di Laboratorio}
% ---------------------------------------------------------------------

Il BST generico con gli algoritmi delle slide: inserimento iterativo con il padre,
ricerca iterativa, minimo e massimo, le tre visite ricorsive, altezza, \cc|trapianta| e
cancellazione nei tre casi, distruttore in post-order. Il \cc|main| prova la
cancellazione di una foglia, di un nodo con un figlio, di uno con due figli e della radice.

\codice{codice/m13/pattern.cpp}
\uscita{codice/m13/output/pattern.txt}

\textbf{Cosa osservare:} dopo ogni cancellazione la visita in-order resta
\textbf{ordinata}: è il modo più rapido per controllare di non aver rotto il BST.
Togliendo la radice 50 (due figli) sale 65, il minimo del sottoalbero destro.

% ---------------------------------------------------------------------
\section{Tips \& Trappole da Esame}
% ---------------------------------------------------------------------

\begin{trappola}[: dimenticare il caso dell'albero vuoto]
\cc|root->val| con \cc|root == nullptr| è un crash. Ogni funzione ricorsiva parte da
\cc|if (p == nullptr) ...|; \cc|min()| e \cc|max()| su un albero vuoto vanno gestiti a parte.
\end{trappola}

\begin{trappola}[: controllare solo i figli]
Verificare che ogni nodo sia maggiore del figlio sinistro e minore del destro
\textbf{non basta} per dire che è un BST: conta tutto il sottoalbero
(Esercizio~13.8).
\end{trappola}

\begin{trappola}[: inserimento ricorsivo senza assegnare]
\cc|inserisci(r->left, v);| senza \cc|r->left = ...| perde il nuovo nodo: la funzione
restituisce il nuovo sottoalbero, e il risultato va \textbf{agganciato}. Lo stesso vale
per \cc|radice = inserisci(radice, v);|.
\end{trappola}

\begin{trappola}[: distruggere in pre-order]
\cc|delete p; distruggi(p->left);| legge un nodo già cancellato. Prima i figli, poi il
padre: post-order.
\end{trappola}

\begin{trappola}[: duplicati]
Decidi una regola e rispettala in \emph{tutte} le funzioni: nelle slide i valori uguali
vanno a \textbf{destra} (\cc|<| a sinistra, \cc|>=| a destra). Se ricerca e
inserimento usano regole diverse, la ricerca non trova i duplicati.
\end{trappola}

\begin{consiglio}[: le funzioni ricorsive ``privata + pubblica'']
Nella classe, la versione ricorsiva lavora su un \cc|Nodo*| (privata); quella pubblica
non ha parametri e la chiama sulla radice: \cc|int altezza() const { return altezza(root); }|.
Così chi usa la classe non vede mai i nodi.
\end{consiglio}

% ---------------------------------------------------------------------
\section{Quiz a Crocette (Prova Scritta)}
% ---------------------------------------------------------------------

\begin{quiz}{Quale visita di un BST restituisce i valori in ordine crescente?}
  \item pre-order
  \item in-order
  \item post-order
  \item per livelli
\end{quiz}

\begin{quiz}{Si inseriscono nell'ordine 50, 30, 70, 20, 40 in un BST vuoto. Qual è la visita pre-order?}
  \item 20 30 40 50 70
  \item 50 30 20 40 70
  \item 20 40 30 70 50
  \item 50 30 70 20 40
\end{quiz}

\begin{quiz}{Qual è la complessità della ricerca in un BST di altezza $h$ con $n$ nodi?}
  \item $O(1)$
  \item $O(h)$, cioè $O(\log n)$ se bilanciato e $O(n)$ nel caso peggiore
  \item sempre $O(\log n)$
  \item sempre $O(n^2)$
\end{quiz}

\begin{quiz}{Inserendo in un BST vuoto i valori 1, 2, 3, \ldots, $n$ in quest'ordine, che albero si ottiene?}
  \item un albero bilanciato di altezza $\log n$
  \item un albero degenere di altezza $n - 1$
  \item un albero con tutte le foglie allo stesso livello
  \item un errore: i valori devono essere inseriti in ordine casuale
\end{quiz}

\begin{quiz}{Dove si trova il minimo di un BST?}
  \item nella radice
  \item nel nodo più a sinistra (si scende sempre a sinistra)
  \item in una foglia qualunque
  \item nel nodo più a destra
\end{quiz}

\begin{quiz}{Cancellando da un BST un nodo con due figli, quale nodo ne prende il posto (nella procedura delle slide)?}
  \item il figlio sinistro
  \item il figlio destro
  \item il successore, cioè il minimo del sottoalbero destro
  \item la radice
\end{quiz}

\begin{quiz}{Quale visita è adatta a distruggere un albero liberando tutti i nodi?}
  \item pre-order
  \item in-order
  \item post-order
  \item nessuna: bisogna usare una coda
\end{quiz}

\begin{quiz}{Quanti nodi può avere al massimo un albero binario di altezza 3?}
  \item 7
  \item 8
  \item 15
  \item 16
\end{quiz}

\begin{quiz}{Dove vengono sempre inseriti i nuovi elementi in un BST?}
  \item come nuova radice
  \item come nuove foglie
  \item a metà dell'albero
  \item nel sottoalbero più basso
\end{quiz}

\begin{quiz}{Un albero binario ha radice 10, figlio sinistro 5 e, come figlio destro di 5, il nodo 12. È un BST?}
  \item sì, perché 12 > 5
  \item no, perché 12 si trova nel sottoalbero sinistro di 10 ma è maggiore di 10
  \item sì, perché ogni nodo ha al più due figli
  \item non si può dire senza conoscere l'altezza
\end{quiz}

\begin{soluzione}[: risposte commentate]
\begin{enumerate}[label=\textbf{\arabic*.}]
  \item \textbf{B}.
  \item \textbf{B}. Radice 50, poi tutto il sottoalbero sinistro (30, 20, 40), poi il destro (70).
        La A è l'in-order, la C il post-order.
  \item \textbf{B}.
  \item \textbf{B}. Ogni nuovo valore va a destra dell'ultimo: è una lista.
  \item \textbf{B}.
  \item \textbf{C}.
  \item \textbf{C}. Prima i figli, poi il padre.
  \item \textbf{C}. $2^{3+1} - 1 = 15$ (1 + 2 + 4 + 8).
  \item \textbf{B}.
  \item \textbf{B}. La proprietà vale per tutto il sottoalbero.
\end{enumerate}
\end{soluzione}

% ---------------------------------------------------------------------
\section{Laboratorio a Step (Simulazione d'Esame)}
% ---------------------------------------------------------------------

\begin{traccia}[: BST Persone (traccia d'esame delle slide)]
Considerando la gerarchia \cc|Persona| $\to$ \cc|Studente|, \cc|Lavoratore|,
\cc|Pensionato|: (a) creare 6 persone (2 per categoria) con nome, cognome ed età casuale
negli intervalli 5--30 (studenti), 20--70 (lavoratori), 65--95 (pensionati); (b) inserirle
in un unico BST con l'età come chiave; (c) scrivere una funzione che calcola e stampa l'età
media per categoria effettuando \textbf{una sola visita} della struttura dati.
\end{traccia}

\begin{step}{un BST chiave--valore}
Scrivi \cc|template <typename K, typename V> class BST| in cui ogni nodo contiene una
\cc|chiave| (che decide la posizione) e un \cc|valore| (il dato associato). Distruttore in
post-order.\\[0.3em]
\textbf{Concetto nuovo:} separare chiave e dato.
\end{step}

\begin{step}{inserimento ricorsivo}
Scrivi l'inserimento in versione ricorsiva (funzione privata che riceve e restituisce un
\cc|Nodo*|) più quella pubblica \cc|insert(K, V)|.\\[0.3em]
\textbf{Concetto nuovo:} ricorsione che restituisce il sottoalbero aggiornato.
\end{step}

\begin{step}{visita con funzione}
Scrivi \cc|visitaInOrder(void (*visita)(K, V))|, che applica la funzione passata a ogni
nodo in ordine di chiave, e \cc|altezza()|.\\[0.3em]
\textbf{Concetto nuovo:} visita generica tramite puntatore a funzione.
\end{step}

\begin{step}{la gerarchia}
Classe astratta \cc|Persona| (nome, cognome, età, \cc|categoria()| virtuale pura) e le
tre derivate.\\[0.3em]
\textbf{Concetto nuovo:} riuso del Modulo 9.
\end{step}

\begin{step}{creazione e inserimento}
Con \cc|srand(111222333)| e una funzione \cc|casuale(a, b)| crea le 6 persone negli
intervalli richiesti e inseriscile in un \cc|BST<int, Persona*>|. Stampale con una visita
in-order (risultano ordinate per età).\\[0.3em]
\textbf{Concetto nuovo:} BST di puntatori a oggetti polimorfici.
\end{step}

\begin{step}{medie con una sola visita}
Con una sola \cc|visitaInOrder| accumula, per ogni categoria (riconosciuta con
\cc|dynamic_cast|), la somma delle età e il numero di persone; poi stampa le tre medie.
Libera la memoria delle persone.\\[0.3em]
\textbf{Concetto nuovo:} calcoli aggregati durante una visita.
\end{step}

\subsection*{Soluzione completa}
\codice{codice/m13/lab.cpp}
\uscita{codice/m13/output/lab.txt}

\begin{consiglio}[: cosa notare nella soluzione]
\begin{itemize}
  \item Il BST contiene \textbf{puntatori}: il suo distruttore libera i nodi, non le
        persone. Le persone le ha create il \cc|main| ed è il \cc|main| a liberarle.
  \item Gli accumulatori sono globali perché la funzione di visita deve avere la firma
        \cc|void(K, V)|. Una soluzione più elegante userebbe un oggetto con i contatori,
        ma questa è più vicina a quanto richiesto in laboratorio.
  \item L'altezza (4) è grande per 6 nodi: le età sono arrivate quasi in ordine crescente,
        e l'albero è sbilanciato.
\end{itemize}
\end{consiglio}

% ---------------------------------------------------------------------
\section{Domande e Risposte Modello per l'Orale}
% ---------------------------------------------------------------------

\begin{orale}[D1. Cos'è un albero binario di ricerca?]
È un albero binario in cui per ogni nodo $x$ tutte le chiavi del sottoalbero sinistro sono
minori della chiave di $x$ e tutte quelle del sottoalbero destro sono maggiori o uguali.
Questa proprietà permette di cercare un valore scendendo dalla radice e scegliendo a ogni
nodo un solo sottoalbero, come nella ricerca binaria. Ricerca, inserimento, cancellazione,
minimo e massimo costano $O(h)$, dove $h$ è l'altezza: $O(\log n)$ se l'albero è bilanciato,
$O(n)$ se è degenere. La visita in-order restituisce le chiavi in ordine crescente.
\end{orale}

\begin{orale}[D2. Quali sono le visite di un albero binario?]
Le tre visite in profondità, definite ricorsivamente, differiscono per il momento in cui si
visita la radice: in pre-order prima la radice, poi il sottoalbero sinistro e il destro; in
in-order il sinistro, la radice, il destro; in post-order il sinistro, il destro e infine la
radice. Hanno tutte costo $\Theta(n)$, perché visitano ogni nodo una volta. La visita per
livelli (in ampiezza) usa invece una coda e visita i nodi livello per livello.
\end{orale}

\begin{orale}[D3. Come si cancella un nodo da un BST?]
Si distinguono tre casi. Se il nodo non ha figli, lo si elimina e il puntatore del padre
diventa nullo. Se ha un solo figlio, il figlio con il suo sottoalbero prende il posto del
nodo. Se ha due figli, si individua il successore, cioè il minimo del sottoalbero destro,
che non ha figlio sinistro: il successore prende il posto del nodo, e se non era figlio
diretto del nodo viene prima sostituito dal proprio figlio destro. Nelle slide gli
spostamenti si fanno con la procedura \emph{Trapianta}, che sostituisce un sottoalbero con
un altro aggiornando il padre. Il costo è $O(h)$.
\end{orale}

\begin{orale}[D4. Perché un BST può degenerare e cosa comporta?]
La forma di un BST dipende dall'ordine di inserimento. Inserendo chiavi già ordinate, ogni
nuovo nodo diventa figlio destro dell'ultimo inserito e l'albero si riduce a una lista di
altezza $n - 1$: tutte le operazioni, che costano $O(h)$, diventano $O(n)$. Un albero
bilanciato ha invece altezza $O(\log n)$. Esistono BST autobilancianti (come gli alberi
AVL o rosso-neri) che mantengono l'altezza logaritmica; nei casi semplici si può
costruire un albero bilanciato inserendo prima il mediano dei valori.
\end{orale}

\begin{orale}[D5. Perché gli algoritmi sugli alberi sono quasi sempre ricorsivi?]
Perché l'albero è una struttura ricorsiva: un albero è vuoto oppure è formato da una radice
e da sottoalberi che sono a loro volta alberi. Un'operazione sull'intero albero si
scompone quindi nell'operazione sulla radice e nella stessa operazione sui sottoalberi, con
caso base l'albero vuoto: visite, altezza, conteggio dei nodi, distruzione. La profondità
della ricorsione è pari all'altezza dell'albero.
\end{orale}

% ---------------------------------------------------------------------
\section{Esercizi al PC}
% ---------------------------------------------------------------------

\begin{esercizio}{Le tre visite}{facile}
Con un nodo semplice (\cc|struct Nodo { int val; Nodo* sx; Nodo* dx; }|) e un inserimento
ricorsivo, costruisci il BST di 50, 30, 70, 20, 40, 60, 80 e stampa le visite pre, in e
post-order. Prima di eseguire, scrivi su carta cosa ti aspetti.
\end{esercizio}

\begin{esercizio}{Contare}{facile}
Scrivi funzioni ricorsive per contare i nodi, contare le foglie e sommare i valori.
Ricava il numero di nodi interni.
\end{esercizio}

\begin{esercizio}{Altezza e livello}{facile}
Scrivi \cc|altezza(r)| (albero vuoto $-1$, foglia 0) e \cc|livello(r, x)|, che restituisce
la profondità del nodo con valore \cc|x| (o $-1$).
\end{esercizio}

\begin{esercizio}{Minimo e massimo}{facile}
Minimo e massimo di un BST, in versione iterativa e ricorsiva.
\end{esercizio}

\begin{esercizio}{Bilanciato o degenere}{medio}
Inserisci i numeri da 1 a 1000 in un BST in ordine crescente, e in un altro inserendo
prima il mediano e poi, ricorsivamente, i mediani delle due metà. Conta i confronti
necessari a cercare 1000 nei due alberi. (Usa un inserimento iterativo: quello ricorsivo
sull'albero degenere farebbe 1000 chiamate annidate.)
\end{esercizio}

\begin{esercizio}{Visita per livelli}{medio}
Stampa i nodi di un albero livello per livello usando una coda di puntatori a nodo, con il
numero di livello all'inizio di ogni riga.
\end{esercizio}

\begin{esercizio}{k-esimo e intervallo}{medio}
Trova il $k$-esimo valore più piccolo con una visita in-order che conta i nodi, e stampa
tutti i valori compresi in $[a, b]$ visitando solo i sottoalberi che possono contenerne.
\end{esercizio}

\begin{esercizio}{È un BST?}{medio}
Scrivi una funzione che verifica se un albero binario è un BST. Prova prima la versione
``ingenua'' (ogni nodo confrontato solo con i figli) sull'albero 10, 5, 20 con 12 figlio
destro di 5, e mostra che sbaglia; poi la versione corretta, che passa a ogni chiamata
l'intervallo $[\text{min}, \text{max})$ ammesso.
\end{esercizio}

\begin{esercizio}{Da array ordinato a BST bilanciato}{difficile}
Costruisci un BST di altezza minima a partire da un array ordinato: il mediano diventa la
radice, e le due metà diventano ricorsivamente i sottoalberi. Verifica con le visite e
l'altezza per un array di 15 elementi.
\end{esercizio}

\begin{esercizio}{Successore e predecessore}{difficile}
Con nodi dotati di puntatore al padre, implementa \cc|successore(x)| e
\cc|predecessore(x)| come nelle slide. Usali per scorrere tutto l'albero in ordine
partendo dal minimo, senza ricorsione.
\end{esercizio}

\subsection*{Soluzioni}
Le soluzioni commentate, con il codice e l'output reale, sono nel PDF separato
\textbf{Soluzioni} (\texttt{soluzioni-light.pdf} / \texttt{soluzioni-dark.pdf}),
allo stesso numero di esercizio. Il codice è anche in \texttt{codice/m13/}.
